% ==============================================================
% Section: Optimization with CES Objectives and Constraints
% ==============================================================

\chapter{Optimization with CES Objectives and Constraints}
\label{sec:ces-constraints}

This section studies optimization when \emph{both} the objective and the constraint have CES structure. Such problems arise naturally in multi-stage production, transformation frontiers, nested allocation decisions, and---as we shall see---intertemporal consumption with non-trivial savings technology.

% --------------------------------------------------------------
\section{The Problem}
\label{subsec:ces-ces-problem}
% --------------------------------------------------------------

\begin{definition}[CES/CES Optimization Problem]
\label{def:ces-ces-problem}
Maximize a CES objective subject to a CES constraint:
\begin{equation}
\label{eq:ces-ces-objective}
\U\opt = \left( \omega_1^{1-q} c_1^{q} + \omega_2^{1-q} c_2^{q} \right)^{1/q}, \qquad q \leq 1,
\end{equation}
subject to
\begin{equation}
\label{eq:ces-ces-constraint}
\Hfun(c_1, c_2) = \left( \eta_1^{1-h} c_1^{h} + \eta_2^{1-h} c_2^{h} \right)^{1/h} = 1, \qquad h \geq 1,
\end{equation}
where $\omega_1 + \omega_2 = 1$, $\eta_1 + \eta_2 = 1$, and at least one of the inequalities $q < 1$ or $h > 1$ is strict.
\end{definition}

The constraint $\Hfun = 1$ defines a \emph{transformation frontier} or \emph{production possibilities frontier} when interpreted as a technology. The curvature of the constraint ($h \geq 1$) opposes the curvature of the objective ($q \leq 1$).

\begin{calloutnote}[When This Structure Arises]
\begin{enumerate}[label=(\roman*), leftmargin=*]
    \item \textbf{Nested production}: Intermediate goods combine via CES to produce final output
    \item \textbf{Transformation frontiers}: Resources transform between uses with diminishing returns
    \item \textbf{Two-stage allocation}: First allocate across categories, then within categories
\end{enumerate}
\end{calloutnote}

% --------------------------------------------------------------
\section{Direct Solution via First-Order Conditions}
\label{subsec:ces-ces-direct}
% --------------------------------------------------------------

The direct approach applies Lagrangian methods to the original problem.

\begin{proposition}[FOC for CES/CES Problem]
\label{prop:ces-ces-foc}
The first-order conditions yield:
\begin{equation}
\label{eq:ces-ces-foc}
\frac{\omega_1^{1-q} c_1^{q-1}}{\omega_2^{1-q} c_2^{q-1}} = \frac{\eta_1^{1-h} c_1^{h-1}}{\eta_2^{1-h} c_2^{h-1}}.
\end{equation}
\end{proposition}

\begin{proof}
The Lagrangian is $\mathcal{L} = \U - \lambda(\Hfun - 1)$. Differentiating with respect to $c_1$ and $c_2$ and taking the ratio eliminates the multiplier $\lambda$ and the common terms from $\U^{1-q}$ and $\Hfun^{1-h}$.
\end{proof}

Rearranging \eqref{eq:ces-ces-foc}:
\begin{equation}
\label{eq:ces-ces-foc-rearranged}
\frac{\omega_1^{1-q}}{\omega_2^{1-q}} \cdot \frac{c_1^{q-h}}{c_2^{q-h}} = \frac{\eta_1^{1-h}}{\eta_2^{1-h}}.
\end{equation}

This can be solved directly, but the expressions become cumbersome. The change-of-coordinates approach below yields cleaner formulas and deeper insight.

% --------------------------------------------------------------
\section{Solution via Change of Coordinates}
\label{subsec:ces-ces-coordinates}
% --------------------------------------------------------------

The key insight is to transform variables so that the CES constraint becomes \emph{linear}. This allows us to apply the canonical formulas from the previous section.

\begin{definition}[Coordinate Transformation]
\label{def:coordinate-transform}
Define new variables:
\begin{equation}
\label{eq:transform-forward}
x_1 = \eta_1^{1-h} c_1^{h}, \qquad x_2 = \eta_2^{1-h} c_2^{h}.
\end{equation}
The inverse transformation is:
\begin{equation}
\label{eq:transform-inverse}
c_1 = \eta_1^{\frac{h-1}{h}} x_1^{1/h}, \qquad c_2 = \eta_2^{\frac{h-1}{h}} x_2^{1/h}.
\end{equation}
\end{definition}

\begin{proposition}[Linearized Constraint]
\label{prop:linear-constraint}
Under the transformation \eqref{eq:transform-forward}, the CES constraint becomes linear:
\begin{equation}
\label{eq:linear-constraint}
\eta_1 x_1 + \eta_2 x_2 = 1.
\end{equation}
\end{proposition}

\begin{proof}
Substituting into $\Hfun^h = 1$:
\[
\eta_1^{1-h} c_1^{h} + \eta_2^{1-h} c_2^{h} = x_1 + x_2 = 1.
\]
Rewriting with the weights explicit: $\eta_1 \cdot (\eta_1^{-1} x_1) + \eta_2 \cdot (\eta_2^{-1} x_2) = 1$, but since we want the constraint in the form of prices times quantities, we note that the effective ``prices'' are $\eta_1$ and $\eta_2$, giving \eqref{eq:linear-constraint}.
\end{proof}

\begin{remark}
The transformation interprets $\eta_i$ as ``prices'' in the new coordinate system. The non-linearity of the original constraint is absorbed into the change of units.
\end{remark}

% --------------------------------------------------------------
\subsection{The Transformed Objective}
\label{subsec:transformed-objective}
% --------------------------------------------------------------

Substituting \eqref{eq:transform-inverse} into the objective:

\begin{align}
\U &= \left( \omega_1^{1-q} c_1^{q} + \omega_2^{1-q} c_2^{q} \right)^{1/q} \notag \\
&= \left( \omega_1^{1-q} \eta_1^{\frac{(h-1)q}{h}} x_1^{q/h} + \omega_2^{1-q} \eta_2^{\frac{(h-1)q}{h}} x_2^{q/h} \right)^{1/q}.
\label{eq:transformed-objective-raw}
\end{align}

\begin{definition}[Effective Power Parameter]
\label{def:effective-power}
The \emph{effective power} in the transformed problem is:
\begin{equation}
\label{eq:psi-def}
\boxed{\psi = \frac{q}{h}}
\end{equation}
This is the ratio of the objective's power to the constraint's power.
\end{definition}

\begin{calloutimportant}[The Effective Elasticity]
The parameter $\psi$ governs the transformed problem. Since $q \leq 1$ and $h \geq 1$:
\begin{itemize}[leftmargin=*]
    \item $\psi < 1$: The objective is ``more concave'' than the constraint is ``convex''
    \item $\psi = 1$: The curvatures match (e.g., Cobb--Douglas objective with Cobb--Douglas constraint)
    \item $\psi > 1$: The constraint curvature dominates (possible only if $q > h$, which requires $q > 1$)
\end{itemize}
The effective elasticity of substitution in the transformed problem is $\tilde{\varepsilon} = 1/(1-\psi)$.
\end{calloutimportant}

% --------------------------------------------------------------
\subsection{Normalizing the Transformed Weights}
\label{subsec:normalize-weights}
% --------------------------------------------------------------

The expression \eqref{eq:transformed-objective-raw} is not yet in canonical form because the coefficients do not sum to one. We proceed in steps.

\textbf{Step 1: Identify the coefficient structure.}

Rewrite \eqref{eq:transformed-objective-raw} as:
\begin{equation}
\U = \left( \alpha_1 x_1^{\psi} + \alpha_2 x_2^{\psi} \right)^{1/q},
\end{equation}
where $\alpha_i = \omega_i^{1-q} \eta_i^{(h-1)q/h} = \omega_i^{1-q} \eta_i^{q - q/h}$.

\textbf{Step 2: Apply a monotone transformation.}

Since we maximize utility, raising to any positive power preserves the optimum. Define:
\begin{equation}
\tilde{\U} = \U^{h} = \left( \alpha_1 x_1^{\psi} + \alpha_2 x_2^{\psi} \right)^{h/q} = \left( \alpha_1 x_1^{\psi} + \alpha_2 x_2^{\psi} \right)^{1/\psi}.
\end{equation}

\textbf{Step 3: Extract the scaling factor.}

To put this in canonical form with weights summing to one, define:
\begin{equation}
\label{eq:kappa-def}
\kappa = \alpha_1^{\frac{1}{1-\psi}} + \alpha_2^{\frac{1}{1-\psi}} = \left( \omega_1^{1-q} \eta_1^{q-\psi} \right)^{\frac{1}{1-\psi}} + \left( \omega_2^{1-q} \eta_2^{q-\psi} \right)^{\frac{1}{1-\psi}}.
\end{equation}

\begin{definition}[Transformed Weights]
\label{def:transformed-weights}
The \emph{transformed weights} are:
\begin{equation}
\label{eq:tilde-omega}
\tilde{\omega}_i = \frac{\alpha_i^{\frac{1}{1-\psi}}}{\kappa} = \frac{\left( \omega_i^{1-q} \eta_i^{q-\psi} \right)^{\frac{1}{1-\psi}}}{\kappa}, \qquad i = 1, 2.
\end{equation}
By construction, $\tilde{\omega}_1 + \tilde{\omega}_2 = 1$.
\end{definition}

\begin{proposition}[Canonical Form of Transformed Problem]
\label{prop:canonical-transformed}
The optimization problem becomes:
\begin{equation}
\label{eq:canonical-transformed}
\max_{\{x_1, x_2\}} \; \tilde{\U} = \kappa^{\frac{1-\psi}{\psi}} \left( \tilde{\omega}_1^{1-\psi} x_1^{\psi} + \tilde{\omega}_2^{1-\psi} x_2^{\psi} \right)^{1/\psi}
\end{equation}
subject to
\begin{equation}
\eta_1 x_1 + \eta_2 x_2 = 1.
\end{equation}
\end{proposition}

Since $\kappa^{(1-\psi)/\psi}$ is a positive constant, maximizing $\tilde{\U}$ is equivalent to maximizing the canonical CES mean:
\begin{equation}
\M(\bx; \tilde{\bom}, \psi) = \left( \tilde{\omega}_1^{1-\psi} x_1^{\psi} + \tilde{\omega}_2^{1-\psi} x_2^{\psi} \right)^{1/\psi}.
\end{equation}

% --------------------------------------------------------------
\subsection{Applying the Canonical Solution}
\label{subsec:apply-canonical}
% --------------------------------------------------------------

We now have a CES objective (with power $\psi$ and weights $\tilde{\bom}$) subject to a linear constraint (with ``prices'' $\eta_i$). The canonical formulas apply directly.

\begin{proposition}[Solution in Transformed Coordinates]
\label{prop:solution-transformed}
The optimal allocation in the transformed space is:
\begin{equation}
\label{eq:x-star}
x_i\opt = \frac{\tilde{\omega}_i \eta_i^{\frac{1}{\psi - 1}}}{\tilde{\omega}_1 \eta_1^{\frac{\psi}{\psi-1}} + \tilde{\omega}_2 \eta_2^{\frac{\psi}{\psi-1}}}, \qquad i = 1, 2.
\end{equation}
The expenditure shares in the transformed space are:
\begin{equation}
\label{eq:shares-transformed}
\eta_i x_i\opt = \frac{\tilde{\omega}_i \eta_i^{\frac{\psi}{\psi - 1}}}{\tilde{\omega}_1 \eta_1^{\frac{\psi}{\psi-1}} + \tilde{\omega}_2 \eta_2^{\frac{\psi}{\psi-1}}}.
\end{equation}
\end{proposition}

\begin{proof}
Apply the canonical CES demand formula with prices $p_i = \eta_i$, weights $\tilde{\omega}_i$, power $\psi$, and budget $B = 1$.
\end{proof}

\begin{callouttip}[Interpretation]
In the transformed coordinate system:
\begin{itemize}[leftmargin=*]
    \item The constraint weights $\eta_i$ act as ``prices''
    \item The transformed weights $\tilde{\omega}_i$ combine both the objective weights $\omega_i$ and the constraint weights $\eta_i$
    \item The effective elasticity $\tilde{\varepsilon} = 1/(1-\psi)$ governs substitution patterns
\end{itemize}
\end{callouttip}

% --------------------------------------------------------------
\subsection{Returning to Original Coordinates}
\label{subsec:return-original}
% --------------------------------------------------------------

The optimal allocation in original coordinates is recovered via \eqref{eq:transform-inverse}:
\begin{equation}
\label{eq:c-star}
c_i\opt = \eta_i^{\frac{h-1}{h}} \left( x_i\opt \right)^{1/h}, \qquad i = 1, 2.
\end{equation}

\begin{proposition}[Optimal Value]
\label{prop:optimal-value}
The optimal utility is:
\begin{equation}
\label{eq:U-star}
\U\opt = \kappa^{\frac{1-\psi}{q}} \cdot \left( \tilde{\Pstar} \right)^{-h/q},
\end{equation}
where $\tilde{\Pstar}$ is the ideal price index in the transformed problem:
\begin{equation}
\tilde{\Pstar} = \left( \tilde{\omega}_1 \eta_1^{\frac{\psi}{\psi-1}} + \tilde{\omega}_2 \eta_2^{\frac{\psi}{\psi-1}} \right)^{\frac{\psi-1}{\psi}}.
\end{equation}
\end{proposition}

\begin{remark}[Decentralization]
\label{rem:decentralization}
The planning problem can be decentralized as a competitive equilibrium. Suppose the consumer owns a firm operating technology $\Hfun$. The firm maximizes revenue $p_1 y_1 + p_2 y_2$ subject to $\Hfun(y_1, y_2) \leq 1$---a Linear/CES problem. The consumer maximizes $\U(c_1, c_2)$ subject to the budget $p_1 c_1 + p_2 c_2 \leq \Pi$, where $\Pi$ denotes firm profits---a CES/Linear problem. At equilibrium prices, markets clear ($c_i = y_i$) and the allocation coincides with the planning solution.

This illustrates the First Welfare Theorem: competitive equilibria are Pareto efficient. It also shows that the \emph{organization} of production---whether centralized or decentralized---is immaterial for efficiency. Prices encode the technological trade-offs ($\eta_i$, $h$) and transmit them to the consumer, who need not know the technology directly. Exercise~\ref{ex:decentralization} asks you to verify this equivalence.
\end{remark}

% --------------------------------------------------------------
\subsection{Alternative Approach: Linearizing the Objective}
\label{subsec:linear-objective}
% --------------------------------------------------------------

The previous approach transformed coordinates to make the \emph{constraint} linear. We now develop the dual approach: transform coordinates to make the \emph{objective} linear while the constraint remains CES.

\begin{definition}[Objective-Linearizing Transformation]
\label{def:obj-linear-transform}
Define new variables:
\begin{equation}
\label{eq:transform-obj-forward}
y_1 = \omega_1^{1-q} c_1^{q}, \qquad y_2 = \omega_2^{1-q} c_2^{q}.
\end{equation}
The inverse transformation is:
\begin{equation}
\label{eq:transform-obj-inverse}
c_1 = \omega_1^{\frac{q-1}{q}} y_1^{1/q}, \qquad c_2 = \omega_2^{\frac{q-1}{q}} y_2^{1/q}.
\end{equation}
\end{definition}

\begin{proposition}[Linearized Objective]
\label{prop:linear-objective}
Under transformation \eqref{eq:transform-obj-forward}, the CES objective becomes linear:
\begin{equation}
\U^q = y_1 + y_2.
\end{equation}
Since $q \leq 1$, maximizing $\U$ is equivalent to maximizing $y_1 + y_2$.
\end{proposition}

\begin{proposition}[Transformed Constraint]
\label{prop:transformed-constraint}
The CES constraint $\Hfun = 1$ becomes:
\begin{equation}
\label{eq:ces-constraint-transformed}
\left( \tilde{\eta}_1^{1-\phi} y_1^{\phi} + \tilde{\eta}_2^{1-\phi} y_2^{\phi} \right)^{1/\phi} = 1,
\end{equation}
where the effective power is:
\begin{equation}
\boxed{\phi = \frac{h}{q}}
\end{equation}
and the transformed constraint weights $\tilde{\eta}_i$ combine $\omega_i$ and $\eta_i$.
\end{proposition}

\begin{calloutimportant}[Duality of Transformations]
The two approaches are related by:
\begin{center}
\begin{tabular}{lcc}
\toprule
\textbf{Approach} & \textbf{Effective Power} & \textbf{Linearizes} \\
\midrule
Constraint-linearizing & $\psi = q/h$ & Constraint \\
Objective-linearizing & $\phi = h/q = 1/\psi$ & Objective \\
\bottomrule
\end{tabular}
\end{center}
The effective powers are reciprocals: $\phi = 1/\psi$. This reflects a fundamental symmetry: what appears as ``curvature in preferences'' from one viewpoint appears as ``curvature in technology'' from the other.
\end{calloutimportant}

\begin{remark}[Which Transformation to Use?]
The constraint-linearizing approach (Section~\ref{subsec:ces-ces-coordinates}) is typically more convenient because it yields a standard consumer problem with linear budget. The objective-linearizing approach is useful when:
\begin{enumerate}[label=(\roman*)]
    \item The constraint has a natural ``cost minimization'' interpretation
    \item We want to emphasize the role of technology curvature
    \item The problem is naturally cast as minimizing a CES cost subject to a target utility level
\end{enumerate}
\end{remark}

% --------------------------------------------------------------
\subsection{Connection to Intertemporal Elasticity of Substitution}
\label{subsec:ies-connection}
% --------------------------------------------------------------

The CES/CES structure provides a unified framework for understanding how preferences and technology interact in dynamic settings.

\begin{definition}[Intertemporal Problem]
\label{def:intertemporal}
Consider a two-period consumption problem:
\begin{equation}
\max_{c_1, c_2} \; \U = \left( \omega_1^{1-q} c_1^{q} + \omega_2^{1-q} c_2^{q} \right)^{1/q}
\end{equation}
subject to the intertemporal budget:
\begin{equation}
\label{eq:intertemporal-budget}
c_1 + \frac{c_2}{R} = W,
\end{equation}
where $R$ is the gross interest rate and $W$ is initial wealth.
\end{definition}

The standard analysis treats \eqref{eq:intertemporal-budget} as a linear constraint with prices $(1, 1/R)$. The intertemporal elasticity of substitution (IES) is:
\begin{equation}
\text{IES} = \frac{1}{1-q}.
\end{equation}

Now consider a generalization where the savings technology has curvature:

\begin{definition}[Non-Linear Savings Technology]
\label{def:nonlinear-savings}
Suppose the transformation between present and future resources is:
\begin{equation}
\label{eq:ces-technology}
\Hfun(c_1, c_2) = \left( \eta_1^{1-h} c_1^{h} + \eta_2^{1-h} c_2^{h} \right)^{1/h} = W,
\end{equation}
with $h \geq 1$. This captures diminishing returns to intertemporal transfer---each additional unit of future consumption costs more in terms of present consumption.
\end{definition}

\begin{proposition}[Effective IES]
\label{prop:effective-ies}
With CES preferences (power $q$) and CES technology (power $h$), the \emph{effective} intertemporal elasticity of substitution is:
\begin{equation}
\boxed{\text{IES}_{\text{eff}} = \frac{1}{1-\psi} = \frac{1}{1 - q/h}}
\end{equation}
\end{proposition}

\begin{calloutnote}[Equivalence Principle]
The following perturbations have \emph{equivalent effects} on the effective IES:
\begin{enumerate}[label=(\roman*)]
    \item \textbf{Lower preference elasticity}: Decreasing $q$ (more concave utility) reduces $\psi = q/h$, lowering the effective IES
    \item \textbf{Higher technology curvature}: Increasing $h$ (more convex constraint) also reduces $\psi$, with the same effect on the effective IES
\end{enumerate}
In other words: \emph{a lower elasticity of substitution in the technology is equivalent to a lower elasticity of substitution in preferences}, and vice versa. The agent cannot distinguish whether reluctance to substitute intertemporally comes from curved preferences or curved technology.
\end{calloutnote}

\begin{remark}[Interpretation]
This equivalence has important implications:
\begin{itemize}[leftmargin=*]
    \item Empirical estimates of IES may conflate preference and technology parameters
    \item Policy interventions that change technology curvature (e.g., progressive taxation of savings) are equivalent to changing the agent's preferences
    \item In calibration, one can normalize either $q$ or $h$ without loss of generality; only the ratio $q/h$ matters for allocations
\end{itemize}
\end{remark}

% --------------------------------------------------------------
\subsection{Exercise: The Dual Transformation}
\label{subsec:exercise-dual}
% --------------------------------------------------------------

\begin{tcolorbox}[
    enhanced, breakable,
    colback=gray!5, colframe=gray!50!black,
    title={\textbf{Exercise: Objective-Linearizing Transformation}},
    fonttitle=\bfseries
]
Consider the CES/CES problem from Definition~\ref{def:ces-ces-problem}.

\textbf{Part (a):} Starting from the transformation $y_i = \omega_i^{1-q} c_i^q$, show that:
\begin{enumerate}[label=(\roman*)]
    \item The objective becomes $\U = (y_1 + y_2)^{1/q}$
    \item The constraint becomes a CES aggregator in $y$ with power $\phi = h/q$
\end{enumerate}

\textbf{Part (b):} Derive the transformed constraint weights $\tilde{\eta}_i$ as functions of $\omega_i$, $\eta_i$, $q$, and $h$.

\textbf{Part (c):} Solve for the optimal $y_i\opt$ using the Linear/CES formulas, then transform back to obtain $c_i\opt$.

\textbf{Part (d):} Verify that your solution matches the one obtained via the constraint-linearizing approach in Section~\ref{subsec:apply-canonical}.

\textbf{Part (e):} Interpret: Under what economic circumstances is this ``dual'' approach more natural?

\medskip
\emph{Hint for (b)}: Substitute $c_i = \omega_i^{(q-1)/q} y_i^{1/q}$ into $\Hfun^h$ and collect terms.
\end{tcolorbox}

\begin{tcolorbox}[
    enhanced, breakable,
    colback=gray!5, colframe=gray!50!black,
    title={\textbf{Exercise: Decentralization}},
    fonttitle=\bfseries,
    label={ex:decentralization}
]
Consider the CES/CES planning problem from Definition~\ref{def:ces-ces-problem}. We decentralize it into a competitive equilibrium with a consumer and a firm.

\textbf{Setup:}
\begin{itemize}[leftmargin=*]
    \item \textbf{Firm:} Maximizes revenue $p_1 y_1 + p_2 y_2$ subject to $\Hfun(y_1, y_2) \leq 1$
    \item \textbf{Consumer:} Owns the firm, receives profits $\Pi = p_1 y_1^\ast + p_2 y_2^\ast$, and maximizes $\U(c_1, c_2)$ subject to $p_1 c_1 + p_2 c_2 \leq \Pi$
    \item \textbf{Market clearing:} $c_i = y_i$ for $i = 1, 2$
\end{itemize}

\textbf{Part (a):} Write down the firm's first-order conditions. Show that optimal supply satisfies:
\begin{equation*}
\frac{p_1}{p_2} = \frac{\eta_1^{1-h} (y_1^\ast)^{h-1}}{\eta_2^{1-h} (y_2^\ast)^{h-1}}.
\end{equation*}

\textbf{Part (b):} Write down the consumer's first-order conditions. Show that optimal demand satisfies:
\begin{equation*}
\frac{p_1}{p_2} = \frac{\omega_1^{1-q} (c_1^\ast)^{q-1}}{\omega_2^{1-q} (c_2^\ast)^{q-1}}.
\end{equation*}

\textbf{Part (c):} Impose market clearing $c_i = y_i$. Show that the equilibrium conditions from (a) and (b) together imply the planning FOC \eqref{eq:ces-ces-foc}.

\textbf{Part (d):} Verify Walras' Law: show that if the goods market clears, the budget constraint holds with equality.

\textbf{Part (e):} Interpret: The consumer never observes the technology parameters $(\eta_1, \eta_2, h)$ directly. How does the price system transmit this information?

\medskip
\emph{This exercise illustrates the First Welfare Theorem: the competitive equilibrium replicates the efficient planning allocation.}
\end{tcolorbox}

% --------------------------------------------------------------
\subsection{Special Cases}
\label{subsec:ces-ces-special}
% --------------------------------------------------------------

\begin{proposition}[Cobb--Douglas Limit]
\label{prop:cd-limit}
When $q \to 0$ and $h \to 1$ such that $\psi = q/h \to 0$:
\begin{enumerate}[label=(\roman*)]
    \item The objective becomes Cobb--Douglas: $\U \to c_1^{\omega_1} c_2^{\omega_2}$
    \item The constraint becomes linear: $\Hfun \to \eta_1 c_1 + \eta_2 c_2$
    \item The solution reduces to $\eta_i c_i\opt = \omega_i$
\end{enumerate}
\end{proposition}

\begin{proposition}[Symmetric Case]
\label{prop:symmetric}
When $\omega_1 = \omega_2 = 1/2$ and $\eta_1 = \eta_2 = 1/2$:
\begin{enumerate}[label=(\roman*)]
    \item The transformed weights are also symmetric: $\tilde{\omega}_1 = \tilde{\omega}_2 = 1/2$
    \item The optimal allocation is $c_1\opt = c_2\opt$
    \item The optimal utility is $\U\opt = 2^{(1-h)/h}$
\end{enumerate}
\end{proposition}

% --------------------------------------------------------------
\subsection{Comparative Statics}
\label{subsec:ces-ces-compar}
% --------------------------------------------------------------

The response of optimal allocations to parameter changes involves three channels:

\begin{enumerate}[label=\textbf{(\arabic*)}]
    \item \textbf{Direct effect}: Changes in $\omega_i$ or $\eta_i$ directly affect the transformed weights $\tilde{\omega}_i$
    
    \item \textbf{Price effect}: Changes in $\eta_i$ affect both the ``prices'' in the linear constraint and the transformed weights
    
    \item \textbf{Curvature effect}: Changes in $q$ or $h$ affect the effective power $\psi = q/h$, which governs substitution patterns
\end{enumerate}

\begin{calloutnote}[The Key Insight]
The effective elasticity is the \emph{ratio} $\psi = q/h$. This means:
\begin{itemize}[leftmargin=*]
    \item Increasing objective concavity ($q \downarrow$) reduces $\psi$, making demand less elastic
    \item Increasing constraint convexity ($h \uparrow$) also reduces $\psi$, with the same effect
    \item The two curvatures interact multiplicatively, not additively
\end{itemize}
\end{calloutnote}

The full comparative statics analysis---including decomposition into income-like and substitution-like effects---is developed in Chapter~\ref{ch:recursive-structure}, where we treat nested CES structures more generally.

% --------------------------------------------------------------
\subsection{Summary}
\label{subsec:ces-ces-summary}
% --------------------------------------------------------------

\begin{calloutimportant}[Key Results for CES/CES Optimization]
\begin{enumerate}[label=(\roman*), leftmargin=*]
    \item The \textbf{change of coordinates} $x_i = \eta_i^{1-h} c_i^h$ converts the CES constraint to a linear constraint with ``prices'' $\eta_i$
    
    \item The \textbf{effective power} is $\psi = q/h$, the ratio of objective power to constraint power
    
    \item The \textbf{transformed weights} $\tilde{\omega}_i$ combine both objective weights $\omega_i$ and constraint weights $\eta_i$
    
    \item The \textbf{canonical solution} applies in the transformed space; invert to recover original allocations
    
    \item The \textbf{scaling factor} $\kappa$ captures the welfare cost of the non-linear constraint structure
\end{enumerate}
\end{calloutimportant}
